\chapter{MỞ ĐẦU}
\label{chap:introduction}

\section{Bối cảnh và động lực}

Bệnh Parkinson là một trong những bệnh thoái hoá thần kinh phổ biến ở người cao tuổi.
Giai đoạn đầu của bệnh thường diễn biến âm thầm, khiến việc phát hiện sớm trở nên khó khăn
nếu chỉ dựa vào các triệu chứng vận động điển hình. Một trong những dấu hiệu sớm và dễ
quan sát của bệnh là thay đổi giọng nói: người bệnh thường xuất hiện run giọng, giọng nhỏ,
khàn hoặc mất điều khiển tần số, biên độ. Các đại lượng âm học thu được từ giọng nói do đó
trở thành một nguồn tín hiệu phi xâm lấn, có tiềm năng hỗ trợ sàng lọc và theo dõi bệnh.

Trong công trình nghiên cứu nổi bật, Little và cộng sự đã thu thập và công bố bộ dữ liệu
giọng nói Parkinson trên kho lưu trữ học máy UCI \cite{little2009,uci_parkinsons}. Bộ dữ liệu
gồm 195 bản ghi, mỗi bản ghi được mô tả bởi 22 đặc trưng giọng nói thuộc các nhóm tần số,
độ run (jitter), độ ồn (shimmer), tỷ lệ nhiễu--giọng và các chỉ số phi tuyến, kèm nhãn
\texttt{status} (0 = khỏe, 1 = Parkinson). Kích thước nhỏ và mất cân bằng nhãn làm cho bộ dữ
liệu này trở thành một tình huống thực tế để rèn nghề khai phá dữ liệu: vừa cần kiểm soát
chất lượng dữ liệu, vừa phải chọn phương pháp đánh giá phù hợp.

Khai phá dữ liệu (data mining) là quá trình khám phá tri thức từ dữ liệu lớn, bao gồm các
bước chuẩn bị, khám phá, tiền xử lý, mô hình hoá và đánh giá. Áp dụng quy trình này vào bài
toán phát hiện Parkinson từ giọng nói không chỉ giúp trả lời được các câu hỏi về chất lượng
dữ liệu và cấu trúc tín hiệu, mà còn cho phép so sánh định lượng các mô hình học máy có giám
sát (phân loại) và không giám sát (phân cụm). Đây là động lực để nhóm chọn đề tài này cho
bài tập lớn môn \textbf{Khai phá dữ liệu}.

\section{Mục tiêu đề tài}

\subsection{Mục tiêu tổng quát}
Xây dựng một quy trình khai phá dữ liệu hoàn chỉnh, tái lập được trên bộ dữ liệu giọng nói
Parkinson của UCI; sử dụng một bộ dữ liệu chuẩn duy nhất, cùng một hạt giống ngẫu nhiên và
cùng một bộ chỉ số đánh giá ở mọi bước, nhằm đảm bảo tính minh bạch và khả năng so sánh.

\subsection{Các câu hỏi nghiên cứu}
\begin{itemize}
    \item \textbf{RQ1 -- Chất lượng dữ liệu:} bộ dữ liệu có đặc điểm gì về kích thước, kiểu
          dữ liệu, nhãn, thang đo, giá trị thiếu/trùng, ngoại lai và độ lệch? Cần xử lý ra sao?
    \item \textbf{RQ2 -- Phân loại có giám sát:} với 22 đặc trưng giọng nói, các mô hình
          KNN, SVM và Random Forest đạt kết quả như thế nào khi đánh giá bằng kiểm định chéo
          phân tầng 10-fold và độ đo F-score?
    \item \textbf{RQ3 -- Phân cụm không giám sát:} cấu trúc cụm tự phát sinh trong không gian
          đặc trưng có phù hợp với ranh giới bệnh lý \texttt{status} hay không (đo bằng ARI)?
\end{itemize}

\subsection{Phạm vi và giới hạn}
Đề tài tập trung vào quy trình khai phá dữ liệu và học máy với bộ dữ liệu công khai có sẵn;
không nhằm mục tiêu chẩn đoán y khoa. Kết quả thực nghiệm được giới hạn ở 195 bản ghi, khoảng
32 chủ thể trong dữ liệu thô, nên chưa khẳng định được hiệu năng trên dân số rộng. Nhóm cũng
không thay đổi, bổ sung hay gán lại nhãn ngoài dữ liệu công bố.

\section{Tổ chức báo cáo}

Báo cáo gồm năm chương:
\begin{itemize}
    \item Chương 1 trình bày bối cảnh, mục tiêu, câu hỏi nghiên cứu và phạm vi.
    \item Chương 2 tổng hợp cơ sở lý thuyết và các phương pháp được sử dụng.
    \item Chương 3 mô tả quy trình, kiến trúc thực hiện, môi trường và cách quản lý công việc.
    \item Chương 4 trình bày kết quả thực nghiệm và đối chiếu.
    \item Chương 5 rút ra kết luận, hạn chế và hướng phát triển.
\end{itemize}

\section{Thành viên và vai trò}

Nhóm gồm 6 thành viên, phân công theo nhiệm vụ chính như sau (chi tiết công việc và
minh chứng nằm ở Chương 3 và Phụ lục A).

\begin{table}[H]
\centering
\caption{Thành viên và vai trò trong đề tài}
\begin{tabular}{|c|p{4.6cm}|p{3.2cm}|p{5.0cm}|}
\hline
\textbf{Mã} & \textbf{Họ và tên} & \textbf{MSSV} & \textbf{Vai trò / nhiệm vụ chính} \\
\hline
M01 & Đặng Hoàng Ân   & 080206014982 & Điều phối, chuẩn bị dữ liệu, kiểm soát gate \\
\hline
M02 & Lê Thiên Lộc    & 079306040024 & Khám phá cấu trúc, kiểu dữ liệu, bản đồ đặc trưng \\
\hline
M03 & Huỳnh Gia Hợp    & 054205008367 & Thống kê mô tả, kiểm tra chất lượng, ngoại lai IQR \\
\hline
M04 & Trần Hoàng Thông & 075206008044 & Tiền xử lý, PCA, trực quan hóa \\
\hline
M05 & Ngô Thị Anh Thư   & 087306008195 & Phân loại, cross-validation, F-score \\
\hline
M06 & Nguyễn Hoàng Bách & 056205005522 & Phân cụm, ARI, tổng hợp kết quả và báo cáo \\
\hline
\end{tabular}
\end{table}
